% !TeX root = ../main.tex
\chapter{Ground Truth Benchmark and Judicial Case Suite}
\label{app:ground_truth_catalog}

\section{Master Ground Truth Dataset Composition}
\label{sec:app_gt_dataset_composition}

\begin{table}[htbp]
    \centering
    \small
    \caption[Master Ground Truth Dataset Partitioning across Evaluation Tiers]{Master Ground Truth Dataset Partitioning across Evaluation Tiers}
    \label{tab:app_gt_composition_summary}
    \begin{tabularx}{\linewidth}{>{\RaggedRight\arraybackslash}p{3.4cm}>{\centering\arraybackslash}p{1.8cm}>{\centering\arraybackslash}p{1.6cm}>{\centering\arraybackslash}p{1.3cm}>{\centering\arraybackslash}p{1.1cm}>{\RaggedRight\arraybackslash}X}
        \toprule
        \textbf{Benchmark Tier} & \shortstack{\textbf{Contra-}\\\textbf{diction}} & \shortstack{\textbf{Entail-}\\\textbf{ment}} & \textbf{Neutral} & \textbf{Total} & \textbf{Primary Evaluation Target} \\
        \midrule
        Tier 1: Authentic Pairs & 35 & 35 & 35 & 105 & Real-world provisions vs.~statutes \\
        Tier 2: Injected Perturbations & 50 & 45 & 45 & 140 & Semantic and numeric fault injection \\
        Tier 3: Judicial Precedents & 45 & 30 & 30 & 105 & Supreme Court preemption rulings \\
        \midrule
        \textbf{Total Benchmark Suite} & \textbf{130} & \textbf{110} & \textbf{110} & \textbf{350} & End-to-end evaluation benchmark \\
        \bottomrule
    \end{tabularx}
\end{table}

\section{Exemplar Benchmark Text Pairs and Gold Annotations}
\label{sec:app_gt_exemplar_pairs}

\begin{table}[htbp]
    \centering
    \scriptsize
    \caption[Exemplar Ground Truth Annotation Pairs]{Exemplar Ground Truth Annotation Pairs}
    \label{tab:app_gt_exemplar_listings}
    \begin{tabularx}{\linewidth}{>{\bfseries}p{1.2cm}>{\RaggedRight\arraybackslash}p{3.0cm}>{\RaggedRight\arraybackslash}p{3.0cm}>{\centering\arraybackslash}p{2.3cm}>{\RaggedRight\arraybackslash}X}
        \toprule
        \textbf{Pair ID} & \textbf{Premise (Controlling Statute)} & \textbf{Hypothesis (Local Provision)} & \textbf{Gold Label} & \textbf{Doctrinal Legal Rationale} \\
        \midrule
        GT-001 & \textbf{RA 7160, \S458(a)(1)(iii):} \textit{The Sangguniang Panlungsod may approve ordinances imposing a fine not exceeding P5,000.00 or imprisonment not exceeding one (1) year.} & \textbf{Draft Ord.~2024-B, \S7:} \textit{Violators shall be penalized with an administrative fine of P10,000.00 or two (2) years imprisonment.} & \textbf{Contradiction} & Direct violation of the statutory penalty ceiling mandated for local government units under the Local Government Code. \\
        \addlinespace[3pt]
        GT-042 & \textbf{RA 9994, \S4(a):} \textit{Senior citizens shall be entitled to the grant of twenty percent (20\%) discount on medicines.} & \textbf{Davao Ord.~0281-12, \S3:} \textit{All commercial pharmacies must provide a 20\% discount on prescription drugs to senior citizens.} & \textbf{Entailment} & Consistent municipal implementation mirroring national statutory entitlements under RA 9994. \\
        \addlinespace[3pt]
        GT-089 & \textbf{RA 8550, \S16:} \textit{The city government shall have jurisdiction over municipal waters as defined in this Code for fishery licensing.} & \textbf{Davao Ord.~0334-12, \S2:} \textit{Heavy motor freight vehicles and cargo trailers are prohibited from secondary residential roads.} & \textbf{Neutral} & Distinct legal domains; traffic management does not intersect municipal fishery licensing jurisdiction. \\
        \bottomrule
    \end{tabularx}
\end{table}

\section{Authentic Supreme Court Jurisprudential Case Suite}
\label{sec:app_sc_case_suite}

\begin{table}[htbp]
    \centering
    \scriptsize
    \caption[Supreme Court Jurisprudential Benchmark Suite]{Supreme Court Jurisprudential Benchmark Suite ($N = 11$ Precedent Controversies)}
    \label{tab:app_supreme_court_benchmark_suite}
    \begin{tabularx}{\linewidth}{>{\RaggedRight\arraybackslash}p{2.5cm}>{\RaggedRight\arraybackslash}p{2.3cm}>{\RaggedRight\arraybackslash}p{2.3cm}>{\RaggedRight\arraybackslash}X>{\centering\arraybackslash}p{2.0cm}}
        \toprule
        \textbf{Case Citation} & \textbf{Challenged Local Act} & \textbf{Controlling Superior Law} & \textbf{Authoritative Judicial Holding} & \textbf{Class} \\
        \midrule
        \textit{Mosqueda v. PBGEA}, G.R.~189185 (2016) & Davao Ord.~0309-07 (Aerial Spraying Ban) & PD 1144; Equal Protection & Annulled: ordinance violated equal protection and usurped national technical regulatory authority. & Contradiction \\
        \addlinespace[2pt]
        \textit{City of Davao v. Aranjuez}, G.R.~137887 (2005) & Davao Ord.~216 (Personnel Benefits) & CSC Omnibus Rules; RA 7160 & Annulled: local allowances exceeded statutory civil service standards and local budget ceilings. & Contradiction \\
        \addlinespace[2pt]
        \textit{Davao City Water Dist. v. Aranjuez}, G.R.~194192 (2015) & Local Board Resolutions & PD 198 (Provincial Water Utilities Act) & Confirmed local water district governance subject to national standards under PD 198. & Entailment \\
        \addlinespace[2pt]
        \textit{People v. Palma}, G.R.~L-44113 (1977) & Davao City Traffic Penal Provision & RA 4136 (Land Transportation Code) & Ordinance penal provisions reconciled with national traffic enforcement jurisdiction. & Neutral \\
        \addlinespace[2pt]
        \textit{SEC v. Davao Light}, G.R.~L-25000 (1969) & Davao Municipal Tax on Utility Franchise & NIRC; Franchise Act & Local tax invalid where state franchise statute reserved exclusive taxing power. & Contradiction \\
        \addlinespace[2pt]
        \textit{City Govt. v. Freight Consolidators}, G.R.~178229 (2009) & Davao City Revenue Code on Forwarding & RA 7160, \S133 (Taxing Limitations) & Invalidated: LGUs are expressly prohibited from levying taxes on goods transported into or through local territory. & Contradiction \\
        \addlinespace[2pt]
        \textit{Davao Int. Port Services v. Abarquez}, G.R.~102132 (1993) & Port CBA vs.~Local Regulation & PD 442 (Labor Code of the Philippines) & Upheld primacy of national statutory labor standards over conflicting local interpretations. & Entailment \\
        \addlinespace[2pt]
        \textit{Apex Mining v. Southeast Mindanao}, G.R.~152613 (2006) & Local Mining Permits (Davao Region) & RA 7942 (Philippine Mining Act) & Local government regulatory scope harmonized within national agency statutory authority. & Neutral \\
        \addlinespace[2pt]
        \textit{Magtajas v. Pryce Properties}, G.R.~111097 (1994) & CDO Ord.~3353 (Casino Prohibition) & PD 1869 (PAGCOR Charter); RA 7160 & Landmark doctrine: local ordinances cannot prohibit activities explicitly authorized by national law. & Contradiction \\
        \addlinespace[2pt]
        \textit{City of Manila v. Laguio}, G.R.~118127 (2005) & Manila Ord.~7783 (Motel/Sauna Prohibition) & Civil Code; 1987 Constitution & Annulled: total prohibition of lawful commercial enterprise exceeded police power. & Contradiction \\
        \addlinespace[2pt]
        \textit{Batangas CATV v. CA}, G.R.~138810 (2004) & Batangas Ord.~3 (CATV Rate Setting) & Act 3846; EO 205; EO 436 (NTC Jurisdiction) & Voided: municipal government lacked delegated authority to regulate cable television transmission fees. & Contradiction \\
        \bottomrule
    \end{tabularx}
\end{table}

\section{Tier 2 Injected-Fault Synthetic Perturbation Schema}
\label{sec:app_tier2_perturbation_schema}

\begin{enumerate}[leftmargin=1.5em, itemsep=4pt]
    \item \textbf{Numeric Penalty Ceiling Violations:} Systematically replacing local fine and confinement terms with values exceeding statutory ceilings in RA~7160 \S458(a)(1)(iii) (\textit{e.g.}, modifying \texttt{P5,000.00} to \texttt{P15,000.00} or 1 year to 3 years).
    \item \textbf{Deontic Modality Inversion:} Transforming mandatory statutory obligations (\texttt{shall}) into discretionary authorizations (\texttt{may}), or converting statutory permissions into categorical municipal prohibitions (\texttt{shall not}).
    \item \textbf{Regulatory Usurpation of National Agency Jurisdiction:} Inserting local licensing prerequisites or unilateral revocation powers over commercial activities explicitly reserved to national bodies (\textit{e.g.}, LTFRB for public utility routes, FDA for pharmaceutical distribution, and NTC for telecommunications).
    \item \textbf{Procedural Appeal Truncation:} Synthetically drafting clauses that eliminate statutory administrative appeal periods or bypass mandatory regional administrative exhaustion under Republic Act No.~7160.
\end{enumerate}
